\documentclass[11pt]{article}
\usepackage[margin=1in]{geometry}
\usepackage{hyperref}

\begin{document}
\noindent
\textbf{To:} Editor-in-Chief\\
\textbf{Journal:} \textit{Sustainable Computing: Informatics and Systems} (Elsevier)\\
\textbf{Submission Type:} Original Research Article\\
\textbf{Date:} November 20, 2026\\

\vspace{1em}
\noindent
\textbf{Subject: Submission of manuscript ``SHIFT-GPU: Shift-Aware, Risk-Calibrated and Physics-Guided Learning for Reliable LLM Inference Hardware Selection''}

\vspace{1em}
\noindent
Dear Editor-in-Chief and Esteemed Members of the Editorial Board,

We are pleased to submit our original research manuscript titled \textbf{``SHIFT-GPU: Shift-Aware, Risk-Calibrated and Physics-Guided Learning for Reliable LLM Inference Hardware Selection''} for consideration as a regular research article in \textit{Sustainable Computing: Informatics and Systems}.

With the massive proliferation of generative artificial intelligence and Large Language Models (LLMs), the carbon footprint, energy consumption, and infrastructure expenditure of continuous autoregressive inference have become central challenges in sustainable computing. To maximize throughput per Watt, modern cloud data centers deploy heterogeneous accelerator clusters comprising diverse GPU architectures across multiple generations (e.g., NVIDIA Hopper, Ada Lovelace, Ampere, Turing, and Volta). However, selecting the optimal GPU for a given LLM and latency Service Level Agreement (SLA) without running exhaustive empirical profiling is an open challenge.

While recent literature (including pioneering work in \textit{Sustainable Computing: Informatics and Systems} by Zhang et al., Kumar et al., and Al-Masri et al.) has explored data-driven configuration optimization and forecasting-based LLM resource management, existing predictive models fail dramatically when deployed on previously unseen hardware architectures. Specifically, when an infrastructure operator introduces an unfamiliar accelerator family, standard tree-based or black-box neural predictors suffer severe covariate shift degradation, experiencing a 300\% surge in prediction error and triggering widespread SLA latency violations.

In this work, we propose \textbf{SHIFT-GPU}, a principled computer science framework uniting:
\begin{enumerate}
    \item \textbf{Physics-Guided Inductive Biases:} We introduce \textbf{PICAN} (\textit{Physics-Informed Cross-Attention Network}), which mathematically couples a deterministic microarchitectural memory-bandwidth roofline lower bound with a bipartite cross-attention representation.
    \item \textbf{Covariate-Shift Conformal Risk Calibration:} We formalize group-conditional conformal prediction adjusted for microarchitectural Mahalanobis distance, providing certified non-asymptotic SLA coverage guarantees even under out-of-distribution hardware shift.
    \item \textbf{Regret-Optimal Active Acquisition (ROAA):} We formulate an active few-shot profiling policy that selects informative hardware calibration measurements ($k \in \{1, 3, 5\}$ shots) to directly minimize operational deployment regret rather than passive statistical variance.
\end{enumerate}

Extensive empirical evaluations across 4,163 real-world production configurations (encompassing 42 open-source LLMs and 8 server-grade GPU types) confirm that SHIFT-GPU reduces out-of-distribution inter-token latency error from 11.28\% to 6.12\%, reduces SLA violation rates from 26.7\% to 2.1\%, and cuts hardware profiling overhead by over 95\% compared to standard sweeps.

This manuscript represents original research that has not been published elsewhere and is not under consideration by any other journal or conference. All authors have approved the submission.

Thank you very much for your time and consideration.

\vspace{1.5em}
\noindent
Sincerely,\\
\textbf{The SHIFT-GPU Research Team}
\end{document}
